\subsection{The original pointwise orbital estimate at rank zero}
\label{rank-zero-pointwise}

The compact terminal case admits a direct argument.  It supplies the
original pointwise theorem at $r=0$, including its quotient measures,
rather than an almost-everywhere substitute.  This specialization does
not assert that the ordinary RH construction has $r=0$.

Put $n=k+1$, $k\ge1$, and $G=U(k,1)$.  Let $L$ be the original
rank-zero local Bessel datum, with discrete-series parameter $\lambda$,
compact parameters $u,u'$, negative vectors $z,z'$ and
$T=-h(z,z)>0$, $T'=-h(z',z')>0$.  Write $J=J_LF$ for its original
compact and unipotent average of a finite matrix coefficient $F$.
Compact Haar measures have mass one; the original unipotent Haar
measure is $\kappa_N\delta_1$, with $0<\kappa_N<\infty$.
For any prescribed exterior section $\mu$, put
\begin{equation}\label{rank-zero-original-norm}
 \mathcal N(F)=\int_{|\zeta|^2>TT'}
       (|\zeta|^2-TT')^{k-1}\|J(\mu(\zeta))\|_{\rm HS}^2\,d\zeta.
\end{equation}
This is exactly the original $\mathrm{jNorm2}$: the source density
is zero for $|\zeta|^2<TT'$, and its boundary circle has planar
measure zero, including when $k=1$ \cite{rank-zero-pointwise-source}.

\begin{theorem}\label{rank-zero-original-pointwise}
Suppose $b\ge5n$, $|\lambda_a|\le b$ for all $a$, and
$|u'_i|\le b$ for all $i$.  For every original point
$\gamma\in\mathrm{Point}(\mathbb C,1)$ with
$\mathcal G(\gamma)\ne0$,
\begin{equation}\label{rank-zero-orbital-bound}
 \mathrm{ptWeight}(k,0,T,T',\gamma)\,
       \mathrm{orbSq}_L(F,\gamma)
 \le b^4\max\{1,\mathcal G(\gamma)^{-1}\}\,\mathcal N(F).
\end{equation}
The left side is understood with the original nonnegative-real
convention.  If $\mathcal G(\gamma)>0$, the stronger bound with
$b^4/\mathcal G(\gamma)$ holds.  Consequently, for $C_g\ge10^6$
and $C_b\ge1$, this proves the original
\texttt{TheoremBPrimePointwise} whenever $D.R.rS=0$.
\end{theorem}

\begin{proof}
We first identify the actual functions and measures.  At $r=0$ the
original unipotent group is $N=\{1\}$, its Lie algebra is zero,
and its pinned Haar construction gives $\mu_N=\kappa_N\delta_1$.
Every point is a single block $(\zeta,1)$ with $\zeta\ne0$.
The original root pattern is empty and the stabilizer
$R_g=N\cap gNg^{-1}$ is trivial.  The root-measure condition
therefore forces $\mu_R=\delta_1$: its support is a singleton and
its pushforward to the zero-dimensional root space has mass one.
The original Weil identity then gives $\nu=\mathrm{pairMeas}$,
by testing indicators of all measurable sets.  Thus every permitted
orbital integral is exactly $J(g)$, retaining the factor
$\kappa_N^2$ in $J$ and $\kappa_N^4$ in its squared norm.  No
volume of a compact stabilizer is inserted in this quotient.

For a representative $g$ one has
$h(gz',z)=\zeta$ and hence, by the reverse Cauchy--Schwarz
inequality for negative vectors, $|\zeta|^2\ge TT'$.
If $\mathcal G=|\zeta|^2-TT'<0$, there are no representatives,
so the defining supremum for $\mathrm{orbSq}$ is zero.
If $\mathcal G>0$, every fiber is one
$K_z\times K_{z'}$ orbit, where $K_z=\operatorname{Stab}(z)$.
Indeed, after sending $z/\sqrt T$ to the last basis vector,
the last coordinate of $gz'$ is fixed and its first $k$
coordinates run over a sphere; $U(k)$ acts transitively on that
sphere.  Two elements with the same image of $z'$ differ on
the right by $K_{z'}$.  The equivariances of $J$ and unitarity
of the compact types show that
\begin{equation}\label{rank-zero-point-identification}
 \mathrm{orbSq}_L(F,\gamma)=\|J(\mu(\zeta))\|_{\rm HS}^2,
 \qquad \mathrm{ptWeight}(k,0,T,T',\gamma)=\mathcal G^{k-1}.
\end{equation}

Choose frames sending $z,z'$ to multiples of the last negative
basis vector and translate $J$ on its two sides.  Haar measure
is unchanged, and the two compact stabilizers become
$K_0=\{\operatorname{diag}(A,1):A\in U(k)\}$.  Call the resulting
function $H$.  Its compact averages have the form
$\kappa_N^2 P_0\sigma(g)Q_0$, with $Q_0$ an intertwiner from
the prescribed right compact type and $P_0^*$ an intertwiner
from the left one.  These maps are obtained by compactly averaging
the original finite coefficient maps; all their images consist of
smooth vectors.  The center of $G$ acts on $\sigma$ by the
character of weight $s=\sum_a\lambda_a$, since the sums of both
compact and full half-root sums are zero.  Thus the finite-dimensional
image of $Q_0$ is invariant under the full maximal compact group
$K=U(k)\times U(1)$.  Its $U(1)$ weight is
\[
                \ell_0=s-\sum_i u'_i.
\]
The same argument applies on the left.  In particular, no bound
on the dimension or multiplicity of a compact type is needed.

Let $Y_1,\ldots,Y_{n^2}$ be a real Hilbert--Schmidt orthonormal
basis of $\mathfrak u(k,1)$ adapted to its compact and noncompact
parts.  The Casimir energy identity, in the normalization
$B(X,Y)=\Re\operatorname{tr}(XY)$, is
\begin{equation}\label{rank-zero-casimir-energy}
 \sum_j\|R_{Y_j}H\|_{L^2(G)}^2
             =(2c_K-c_G)\|H\|_{L^2(G)}^2,
\end{equation}
where
\[
 c_G=\sum_a\lambda_a^2-\|\rho_n\|^2,\qquad
 c_K=\sum_i(u'_i)^2-\|\rho_k\|^2+\ell_0^2.
\]
To check the sign, $\Omega_G=\sum_{\mathfrak p}Y^2-
\sum_{\mathfrak k}Y^2$ and $\Omega_K=-\sum_{\mathfrak k}Y^2$,
so $-\sum_jY_j^2=2\Omega_K-\Omega_G$.  Moving positive root
matrix units past negative ones gives the displayed eigenvalues;
the compact highest weight is $u'-\rho_k$.  Schur orthogonality
makes these finite matrix coefficients and their derivatives
square integrable.  Equivalently one may apply the skew-adjoint
derived operators in the right regular representation, so no
boundary term or formal integration by parts is assumed.

Now $|\ell_0|\le(n+k)b\le2nb$, $c_K\le5n^2b^2$ and
$c_G\ge-n^3$.  Since $b\ge5n$,
\begin{equation}\label{rank-zero-energy-bound}
                   2c_K-c_G\le11n^2b^2\le b^4.
\end{equation}
Let $a_t$ be the positive boost on the first positive and last
negative coordinates, with entries $\cosh t,\sinh t$, and put
$\Phi(t)=H(a_t)$ for $t>0$.  The compact equivariances and central
unitarity imply that its Hilbert--Schmidt norm is the norm of
$J(\mu(\zeta))$ whenever $|\zeta|=\sqrt{TT'}\cosh t$.
The exact terminal integration formula, or its direct
hyperboloid change of variables, has radial density
\[
 w(t)=(\sinh t)^{2k-1}\cosh t,
 \qquad M=\int_0^\infty w(t)\|\Phi(t)\|_{\rm HS}^2\,dt.
\]
In particular the angular integral and planar Jacobian give
\begin{equation}\label{rank-zero-radial-normalization}
                        \mathcal N(F)=2\pi(TT')^k M.
\end{equation}
For clarity the associated Haar formula on $G$ may have an
arbitrary positive constant; it is the same constant in both
integrals in the following energy comparison and cancels.
The boost generator $X$ has $\|X\|_{\rm HS}^2=2$.
On $k_1a_tk_2$, its radial derivative is the right derivative
in $\operatorname{Ad}(k_2^{-1})X$.  Conjugation by $K$ preserves
the Hilbert--Schmidt norm, and the compact factors act unitarily
on $H$ and $\Phi'$.  Cauchy--Schwarz, the Haar formula and
\eqref{rank-zero-energy-bound} therefore yield
\begin{equation}\label{rank-zero-radial-energy}
 E:=\int_0^\infty w(t)\|\Phi'(t)\|_{\rm HS}^2\,dt
                                  \le2b^4 M.
\end{equation}

It remains to turn this actual derivative estimate into point
evaluation.  Fix $t>0$ and put
$\delta=\tanh t/(8k)$ and $I=[t-\delta,t+\delta]$.
Here $\delta\le t/8$ and $s\ge7t/8$ on $I$.  Concavity of
$\tanh$ gives $\coth s\le(8/7)\coth t$ there.  Hence
\[
 (\log w)'(s)=(2k-1)\coth s+\tanh s\le2k\coth s,
 \qquad |\log w(s)-\log w(t)|\le2/7<1.
\]
The fundamental theorem of calculus, Cauchy--Schwarz and
averaging over $s\in I$ give the Hilbert-valued estimate
\[
 \|\Phi(t)\|^2\le\delta^{-1}\int_I\|\Phi(s)\|^2\,ds
                       +2\delta\int_I\|\Phi'(s)\|^2\,ds.
\]
Using $\mathrm e^{-1}w(t)\le w(s)\le\mathrm e w(t)$ and
\eqref{rank-zero-radial-energy}, one obtains
\[
 w(t)\|\Phi(t)\|^2
       \le\mathrm e\left(8k\coth t+\frac{b^4}{2k}\tanh t\right)M.
\]
Multiply by $(TT')^{k-1}/(\sinh t\cosh t)$ and use
\eqref{rank-zero-radial-normalization}.  Since
$\mathcal G=TT'\sinh^2t$, this gives
\[
 \mathcal G^{k-1}\|J(\mu(\zeta))\|^2
 \le \frac{\mathrm e}{2\pi}
           \left(8k+\frac{b^4}{2k}\right)
            \frac{\mathcal N(F)}{\mathcal G}
 \le \frac{b^4}{\mathcal G}\mathcal N(F).
\]
The last estimate uses $\mathrm e/(2\pi)<1/2$, $k\ge1$, $k\le b$ and
$b\ge2$.  All inequalities are homogeneous; when $M=0$ the
same argument gives pointwise vanishing, without division by $M$.

Finally, the original $\mathrm{LocOk}$ with $\lambda_0=\lambda$
has $|\lambda_a|,|u'_i|\le b$.  Its noncompact gap, at $r=0$,
is at least $C_g2^{9/2}(3n^2+n^3)$.  Both signs occur because
the underlying form has signature $(k,1)$; thus this gap is at
most $2b$, implying $b\ge5n$.  For $m=2$ and $C_b\ge1$,
\[
 b^4\max(1,\mathcal G^{-1})
 \le b^{16C_b}\max(1,\mathcal G^{-1})^{4C_b},
\]
which is exactly the original $\mathrm{BpIneq}$ budget.
This proves its universally quantified pointwise interface,
including all original representatives and permitted measures.
\end{proof}

The theorem closes the entire rank-zero local pointwise supplier.
For $r>0$, partial-head kernels still involve unipotent quotient
measures, root stabilizers and their density and gauge factors.
The compact Casimir argument above does not identify those kernels
with matrix coefficients of a fixed $U(k,1)$ discrete series.
That identification or a substitute uniform kernel estimate remains
necessary before this terminal argument can be iterated.
